
\section{Introduction}
Intersections are critical nodes in road networks and account for a substantial share of traffic delay, energy consumption, and road accidents \parencite{farina2026v2n, el2026distributed,bui2017real}. Traditional intersection management relies heavily on signalized control, which regulates conflicting vehicle movements through traffic light cycles. Considerable research has aimed at improving the cycle logic, ranging from actuated and adaptive timing to reinforcement-based-signal control, resulting in moderate reductions in delay and energy consumption \parencite{el2026distributed, 10433940}. However, approaches remain constrained by the signalized architecture, which is not capable of handling stochastic arrivals of vehicles, leading to idle waiting time and frequent stop-and-go behavior, especially during congestion \parencite{el2026distributed}. This stop-and-go behavior is directly relevant for energy consumption of electric vehicles, since frequent acceleration and deceleration increase the power demand of electric vehicles \parencite{kamal2014vehicle,galvin2017energy}.
\\\\
The emergence of Autonomous Vehicles (AVs) offers an opportunity to address these inefficiencies, since AVs are predictable and can adapt easily to dynamic conditions, they are capable of outperforming human drivers at coordinating intersection crossings, presenting opportunities for further delay reductions and throughput increments \parencite{soltani2025communication}. These opportunities have resulted in more research on distributed intersection management strategies that either replace or adjust traditional intersection management strategies. 
\\\\
The largest share of this research is based on Vehicle-to-Vehicle (V2V) , Vehicle-to-Network (V2N), or Vehicle-to-Infrastructure (V2I) communication, which allow vehicles to share their position, speed, and negotiate a crossing sequence. \textcite{el2026distributed} optimize vehicle crossing sequences through virtual platooning based on V2V communication, achieving increased intersection throughput and reduced delay, compared to traditional intersection management strategies. In contrast, \textcite{hasibur2021application} apply a V2I and V2V approach in which traffic signals exchange real-time phase and cycle timing information with connected autonomous vehicles, enabling vehicles to adjust their current speed to pass through the intersection more efficiently and the traffic light to dynamically adjust cycle timing. Results showed a decrease in delay and improved traffic safety. On the other hand, \textcite{bifulco2022decentralized} propose a novel longitudinal control strategy that combines on-board range sensors with V2X communication, through which vehicles exchange information to estimate each other's time to the intersection. Results show a decrease in delay time, improvements in terms of safety, and increased throughput. Although, research shows promising results in terms of throughput, delay, and safety, their reliance on communication introduces practical barriers, including the need for standardized inter-vehicle hardware, and susceptibility to delays in communication \parencite{soltani2025communication,el2026distributed}. Furthermore, V2I approaches require new infrastructure at intersections, which introduce substantial costs and single points of failure \parencite{soltani2025communication}.   
\\\\
To avoid these barriers, a smaller body of research investigates communication-free approaches, in which each vehicle decides how to traverse the intersection solely on its own onboard sensors. \textcite{song2025control} show that their adaptive barrier control function is able to safely navigate an intersection, without inter-vehicle communication.  In addition, the communication-free distributed control algorithm of \textcite{soltani2025communication}, achieved throughput comparable to an actuated signalized intersection, while reducing the average delay, using only the on board sensors to measure distances and vehicle speeds.
The results make communication-free control attractive for future-deployment, as it removes dependency on communication infrastructure.
\\\\
Despite the promising results of communication-free control, development is still in its early stages and therefore the following literature gaps exist. First, most studies evaluate intersection management strategies in simplified environments, where vehicles are modeled as points with simplified vehicle dynamics, rather than physics-based simulators which use realistic vehicle dynamics. Second, most studies examine delay time, throughput, fuel consumption and safety as the relevant performance indicators, but the effect on vehicle energy consumption is rarely studied, despite being a global concern \parencite{hadjigeorgiou2022real}.
\\\\
Therefore this thesis addresses these gaps by evaluating a communication-free distributed barrier control algorithm, against a signalized intersection, which follows a Fully Actuated System (FAS), modeled based on Dutch standards \parencite{wilson2014handboek}. The barrier control algorithm utilizes zero communication between AVs, instead each  AV independently estimates its own distance and the distance of the conflicting AV to the  point of conflict, the intersection. Based on the measurements, the AV adjusts its speed to ensure both AVs do not arrive simultaneously at the conflict point. The barrier and FAS controller are modeled, simulated, and compared in CARLA, an open-source, physics-based autonomous driving simulator \parencite{dosovitskiy2017carla}, under identical conditions. Four different traffic densities are simulated across thirty simulation runs each, where each run varied in initial spawn conditions. By evaluating both control strategies in the exact same simulation environment and only varying the intersection management strategy, the difference in performance can be attributed solely to the control strategy, which is evaluated on four Key Performance Indicators (KPIs): throughput, delay time, energy consumption and safety violations.       
\\\\
Contributions to current research, are the evaluation of the barrier control algorithm in a physics-based simulator with realistic vehicle dynamics, providing a more realistic evaluation than the simplified setups used in prior research. In addition, the evaluation focuses on vehicle energy consumption instead of fuel consumption, which is rarely studied in existing literature. 
\\\\
The objective of this thesis is to evaluate the communication-free distributed barrier control algorithm against the FAS controller in terms of throughput, delay time, energy consumption, and safety violations under  varying initial conditions,  using a single-intersection in CARLA simulator. This leads to the following research question:\\
\textit{What is the performance of a barrier control algorithm for autonomous vehicles compared to a Fully Actuated System controller at a single-intersection in terms of throughput, delay time, energy consumption, and safety, based on multiple CARLA simulations under varying initial conditions?}\\
This research question is divided into the following sub-questions:
\begin{enumerate}
    \item Which system and control parameters are required for safe and stable operation of the distributed control algorithm in CARLA?
    \item What differences in throughput, delay time, energy consumption, and safety are observed between both control strategies under identical conditions?
    \item What is the effect of varying initial conditions on the performance of both control strategies?
\end{enumerate}
The remainder of this report is structured as follows: Section~\ref{System Overview} describes the barrier controller, and the FAS controller. Section~\ref{Methodology} describes the methodology, including research approach, simulation environment, experimental design and parameters, and key performance indicators, Section~\ref{Results} presents the results, which are discussed in section~\ref{Discussion}, which is then followed by the conclusion in section~\ref{Conclusion}.  